\newpage
\section{UVX Guide}

Using Python's toolchain runner (\textbf{uvx}) alongside Node gives you access to a completely different class of system automation tools. Running servers via \lstinline[style=inline]|uvx| provides lightning-fast local execution directly from registries.

\subsection{Official File System and Git Servers}
These servers give the agent native capabilities on the local machine without spawning blind bash wrappers.

\begin{itemize}
    \item \textbf{Git History} (\lstinline[style=inline]|mcp-server-git|): Gives the agent native capabilities to read branch histories, calculate diffs, stage files, and check commits.
    \item \textbf{Local Filesystem} (\lstinline[style=inline]|mcp-server-filesystem|): Opens up file I/O operations (reading directories, verifying structures, changing source code) across any system root you expose to it.
    \item \textbf{Fetch Markdown} (\lstinline[style=inline]|mcp-server-fetch|): A Python alternative to the npx fetch server.
\end{itemize}

\paragraph{Universal UVX Configuration Block.}
When configuring these in your runtime hub (e.g., \lstinline[style=inline]|settings.json| or \lstinline[style=inline]|mcp_config.json|), you can pass arguments to specify allowed directories or execution paths. This can be combined with NPX servers in the same configuration.

\begin{lstlisting}[style=json]
{
  "mcpServers": {
    "git-history": {
      "command": "uvx",
      "args": ["mcp-server-git"]
    },
    "local-filesystem": {
      "command": "uvx",
      "args": ["mcp-server-filesystem", "/"]
    },
    "fetch-markdown": {
      "command": "uvx",
      "args": ["mcp-server-fetch"]
    }
  }
}
\end{lstlisting}
